\section{Experiments}
\label{sec:experiments}

\noindent\textbf{Status of this section.}
The evaluation harness accompanying this paper lives in
\texttt{evaluation\_cognitive/} of the public repository and is currently
released as a skeleton: configuration schemas, dataset-adapter interfaces,
and a reproducibility checklist. The CognitiveBench data generator and the
end-to-end ablation runner are explicitly marked as work-in-progress in
\texttt{evaluation\_cognitive/README.md}. Consistent with this
artifact-first stance, this section describes the research questions, the
evaluation protocol we commit to, and the ablation matrix we have wired up
in the configs, but intentionally does \emph{not} quote headline numeric
results whose reproduction currently depends on an unchecked box in the
checklist. Numeric tables will be added once the checklist is fully
ticked; the corresponding tables and figures are referenced here by label
but deferred. The separate upstream benchmark harness in
\texttt{evaluation\_mem0\_original/} is unchanged from
\cite{mem0} and is \emph{not} used to produce any result in this paper.

We organise the evaluation around three research questions:
\begin{itemize}
    \item \textbf{RQ1}: How does cognitive-inspired memory management affect long-term retention and noise reduction compared to static baselines?
    \item \textbf{RQ2}: What is the individual contribution of each cognitive module (emotion weighting, sleep consolidation, and the top-$k$-weighted meta-learner)?
    \item \textbf{RQ3}: Does our system generalise to real-world long-context benchmarks?
\end{itemize}

\subsection{Experimental Setup}
\label{subsec:setup}

\paragraph{Datasets.} We plan to evaluate on two complementary benchmarks:
\begin{enumerate}
    \item \textbf{CognitiveBench}: an in-repository long-context dialogue simulator generating multi-turn conversations with injected factual anchors, emotional events, recurring patterns, and distractor turns. The generator's interface and seed-manifest schema are specified in \texttt{evaluation\_cognitive/generators/cognitivebench.py} and \texttt{evaluation\_cognitive/configs/cognitivebench\_seeds.yaml}; details and the leak-check protocol are in Appendix~\ref{app:cognitivebench}. The generator itself currently raises \texttt{NotImplementedError}; freezing the seed manifest and implementing the generator is the precondition for any CognitiveBench-derived number reported in a later version of this paper.
    \item \textbf{LoCoMo}~\cite{locomo}: we adapt our system to LoCoMo's evaluation protocol via the \texttt{LoCoMoAdapter} interface in \texttt{evaluation\_cognitive/adapters/locomo.py}. The adapter is interface-only in this release; its \texttt{describe\_metric} method is the single source of truth for how we map LoCoMo's answer attribution to our retention-conditioned retrieval metric, and must be implemented before any LoCoMo number is reported.
\end{enumerate}

\paragraph{Baselines.} The planned comparison set is
(1)~\textbf{Vanilla RAG} with fixed window truncation,
(2)~\textbf{Mem0 Original}~\cite{mem0} without cognitive extensions,
(3)~\textbf{Generative Agents}~\cite{generative_agents} with hourly reflection,
(4)~\textbf{MemGPT}~\cite{memgpt} with hierarchical paging, and
(5)~\textbf{Zep}~\cite{zep} with time-decay summarisation.

\paragraph{Evaluation metrics.} We report:
\begin{itemize}
    \item \textbf{Retention@K}: Percentage of gold turns whose id appears in the top-$K$ retrieved memories for a held-out query.
    \item \textbf{Noise Ratio}: Proportion of distractor / non-gold turns in the top-$K$ retrieved set.
    \item \textbf{Token Efficiency}: Reduction in context tokens consumed during retrieval, relative to Vanilla RAG.
    \item \textbf{Latency Overhead}: Additional inference time introduced by cognitive operations.
\end{itemize}
The precise mapping from LoCoMo's answer-attribution label to Retention@K
is pinned in \texttt{LoCoMoAdapter.describe\_metric}; the three candidate
definitions flagged in our design notes
(answer-accuracy, retrieval-faithfulness, attribution-correctness) must not
be mixed silently.

\paragraph{Evaluation protocol on LoCoMo.} We inject historical dialogue turns sequentially into \textsc{Mem0-Cognitive}'s memory store, perform retrieval only at query time (not per turn) to match LoCoMo's setting, and apply the \texttt{LoCoMoAdapter} metric described above.

\paragraph{Ablation matrix.} The four cells we commit to evaluating are materialised as four YAML files in \texttt{evaluation\_cognitive/configs/}: \texttt{ablation\_full.yaml} (all three modules on), \texttt{ablation\_no\_emotion.yaml}, \texttt{ablation\_no\_consolidation.yaml}, and \texttt{ablation\_no\_meta\_learner.yaml}. The ablation runner at \texttt{evaluation\_cognitive/scripts/run\_ablation.py} is currently an interface-only stub; when implemented it will consume any of these configs and write per-turn outputs under \texttt{evaluation\_cognitive/expected\_outputs/}.

\subsection{Main Results: Long-Term Dialogue Simulation}
\label{subsec:main_results}

Table~\ref{tab:main_results} will report the primary comparison on
CognitiveBench once the generator is frozen and
\texttt{evaluation\_cognitive/scripts/run\_ablation.py} has produced a
stable release under
\texttt{evaluation\_cognitive/expected\_outputs/ablation\_full/}.
In the current artifact release, this table is intentionally left empty;
quoting numeric values here without a reproducible source would violate
the checklist in \texttt{evaluation\_cognitive/README.md}.

\begin{table}[t]
\centering
\small
\begin{tabular}{lcccc}
\toprule
\textbf{System} & \textbf{Retention@10} & \textbf{Noise Ratio} & \textbf{Token Savings} & \textbf{Latency (ms)} \\
\midrule
Vanilla RAG & --- & --- & --- & --- \\
Mem0 Original & --- & --- & --- & --- \\
Generative Agents & --- & --- & --- & --- \\
MemGPT & --- & --- & --- & --- \\
Zep & --- & --- & --- & --- \\
\midrule
\textbf{Mem0-Cognitive (Ours)} & --- & --- & --- & --- \\
\bottomrule
\end{tabular}
\caption{Main results on CognitiveBench. Table intentionally empty in this release; values will be populated from \texttt{evaluation\_cognitive/expected\_outputs/ablation\_full/} once the generator and runner checklist in \texttt{evaluation\_cognitive/README.md} is fully ticked.}
\label{tab:main_results}
\end{table}

\subsection{Ablation Study: Module Contributions}
\label{subsec:ablation}

To isolate the contribution of each cognitive component, we commit to a
factorial ablation over the three modules described in
Section~\ref{sec:methodology}: emotion weighting, sleep consolidation, and
the top-$k$-weighted meta-learner (note that the last module is a
heuristic, not GP-based Bayesian Optimisation; see
\S\ref{sec:methodology}). The four configurations that correspond to the
paper's headline ablation rows (baseline, $-$emotion, $-$consolidation,
$-$meta-learner, full) are wired up as config files under
\texttt{evaluation\_cognitive/configs/}; the remaining cells in
Table~\ref{tab:ablation} will be added when the runner is implemented.

\begin{table}[t]
\centering
\small
\begin{tabular}{ccc|cc}
\toprule
\textbf{Emotion} & \textbf{Consolidation} & \textbf{Meta-Learner} & \textbf{Retention@10} & \textbf{Noise Ratio} \\
\midrule
\xmark & \xmark & \xmark & --- & --- \\
\cmark & \xmark & \xmark & --- & --- \\
\xmark & \cmark & \xmark & --- & --- \\
\xmark & \xmark & \cmark & --- & --- \\
\cmark & \cmark & \xmark & --- & --- \\
\cmark & \xmark & \cmark & --- & --- \\
\xmark & \cmark & \cmark & --- & --- \\
\midrule
\textbf{\cmark} & \textbf{\cmark} & \textbf{\cmark} & --- & --- \\
\bottomrule
\end{tabular}
\caption{Factorial ablation matrix on CognitiveBench. Cells intentionally empty in this release; see \texttt{evaluation\_cognitive/configs/ablation\_*.yaml} for the exact runner configuration of each row.}
\label{tab:ablation}
\end{table}

\subsection{Memory Growth Dynamics}
\label{subsec:memory_growth}

Figure~\ref{fig:memory_growth} will visualise the evolution of stored
memory items as dialogue length grows. As with the tables above, the
figure is a forward reference in this release: the underlying data will
be produced by \texttt{evaluation\_cognitive/scripts/run\_ablation.py}
once implemented, and only then will the plot be generated and linked.
The qualitative prediction is that consolidation induces a stable
equilibrium in store size, while naive accumulation grows linearly and
fixed-window truncation loses long-range dependencies.

\subsection{Real-World Validation on LoCoMo}
\label{subsec:locomo_results}

To probe generalisation beyond our synthetic benchmark, we plan to
evaluate on LoCoMo's test set through
\texttt{evaluation\_cognitive/adapters/locomo.py}. As with
CognitiveBench, Table~\ref{tab:locomo} is intentionally empty in this
release; numeric entries will only appear once the adapter's
\texttt{describe\_metric} is pinned and the runner has emitted a stable
output directory.

\begin{table}[h]
\centering
\small
\begin{tabular}{lc}
\toprule
\textbf{System} & \textbf{Accuracy (\%)} \\
\midrule
Vanilla RAG & --- \\
Mem0 Original & --- \\
Generative Agents & --- \\
MemGPT & --- \\
\midrule
\textbf{Mem0-Cognitive (Ours)} & --- \\
\bottomrule
\end{tabular}
\caption{Accuracy on LoCoMo. Values will be populated from
\texttt{evaluation\_cognitive/expected\_outputs/locomo/} once the
\texttt{LoCoMoAdapter} is implemented and its metric definition is
pinned.}
\label{tab:locomo}
\end{table}

\subsection{Qualitative Analysis: Case Studies}
\label{subsec:case_studies}

When the runner is implemented, two illustrative case studies will be
emitted directly from logged consolidation and retrieval traces: (i) an
emotion-weighted retention case in which a rare but high-emotion
fact persists despite low access frequency, and (ii) a sleep-consolidation
case in which multiple related low-level episodes are abstracted into a
single semantic memory. The corresponding dialogue transcripts and audit
logs will be checked in under
\texttt{evaluation\_cognitive/expected\_outputs/} and cross-referenced
from Appendix~\ref{app:case_studies}.

\subsection{Discussion: Expected Ablation Profile}
\label{subsec:discussion}

We outline the ablation profile we expect, and which the empty cells in
Table~\ref{tab:ablation} are designed to refute or confirm. The working
hypothesis, consistent with the mechanism analysis in
\S\ref{sec:methodology}, is that sleep consolidation is the dominant
contributor to noise reduction, that emotion weighting primarily affects
tail retention of high-salience events, and that the top-$k$-weighted
meta-learner yields modest but monotone per-user improvements on
longer horizons. We deliberately avoid committing to numeric values here
until the runner produces them; a future version of this paper will
either confirm or falsify this profile based on the checked-in outputs.
